\documentclass[../../../../main.tex]{subfiles}
\begin{document}

$xy$ 平面上の長さ2の線分 $AB$ を直径とする半円を $D$ とする．
半円 $D$ の内部(周を含まない)の一点を $P$ とする．
$A$ と $P$ を通る直線と半円 $D$ の円弧の部分との交点を $Q$ とし，
$B$ と $P$ を通る直線と半円 $D$ の円弧の部分との交点を $R$ とする．
五角形 $ARPQB$ の面積を $S$ とおく．

\begin{figure}[htbp]
  \centering
  \begin{tikzpicture}[scale=2, >=stealth]
    \coordinate (A) at (-1, 0);
    \coordinate (B) at (1, 0);
    \coordinate (O) at (0, 0);

    \draw[thick] (A) -- (B);
    \draw[thick] (1,0) arc (0:180:1);

    \coordinate (R) at (135:1);
    \coordinate (Q) at (40:1);

    \coordinate (P) at (intersection of A--Q and B--R);

    \draw (A) -- (Q);
    \draw (B) -- (R);
    \draw (A) -- (R);
    \draw (B) -- (Q);

    \draw (P) +(-40:0.25) arc (-40:-140:0.25);
    \node at ($(P) + (-90:0.4)$) {$\theta$};

    \fill (P) circle (1pt) node[above] {$P$};
    \fill (R) circle (1pt) node[above left] {$R$};
    \fill (Q) circle (1pt) node[above right] {$Q$};

    \node[below left] at (A) {$A$};
    \node[below right] at (B) {$B$};
  \end{tikzpicture}
\end{figure}

\begin{enumerate}
  \item $\angle APB$ を一定に保ったまま点 $P$ が半円 $D$ の内部を動くとき，
  $S$ のとる値の範囲を，$\angle APB=\theta$ を使って表せ．
  \item 点 $P$ が，半円 $D$ の内部を自由に動くとき，$S$ のとる値の範囲を求めよ．
\end{enumerate}

\end{document}
